%%
% 第二章：Skill 的形式化建模与 SkillRank 算法设计
%%

\chapter{Skill 的形式化建模与 SkillRank 算法设计}
\label{cha:modeling}

\section{Agent Skill 的结构分析}
\label{sec:skill_structure}

在现有 Agent 框架（如 Claude Code、Cursor、Manus）的 Skill 生态中，Skill 的外在形式多种多样：有以 Markdown 格式编写的操作规范、以 YAML 描述的工具调用配置、以 JSON 定义的工作流模板，也有直接注入系统提示（System Prompt）的纯文本片段。尽管形式各异，这些 Skill 在功能上都可以被统一归纳为一个二元组结构：

\begin{equation}
    \label{eq:skill_def}
    s = \langle \mathit{Description},\ \mathit{Content} \rangle
\end{equation}

其中，$\mathit{Description}$ 是 Skill 的功能声明字段，描述该 Skill 能做什么、适用于哪类任务、在何种条件下应当被激活。$\mathit{Content}$ 是 Skill 的实质内容，即 Skill 被触发后实际写入 Agent 上下文窗口的文本——它可能是一段操作规范、一组最佳实践、一个领域知识摘要，或者是一套约束条件与输出格式要求。

这两个字段在 Skill 的使用生命周期中承担着完全不同的职能，这一区分是本文整个分析体系的基础。当 Agent 接收到一个任务时，它首先扫描可用 Skill 的 Description 字段，判断哪些 Skill 与当前任务在语义上相关，并决定是否将其激活。这一过程可以理解为一个\textbf{语义定位}操作：Agent 在由所有 Skill Description 定义的语义空间中，找到与任务最近邻的若干个 Skill。一旦某个 Skill 被选中，其 Content 便被复制到 Agent 的上下文中，从此刻起直接参与 Agent 的推理过程，影响其最终输出的质量与准确性。这一过程是一个\textbf{质量贡献}操作：高质量的 Content 能够使 Agent 在该领域内展现出接近专家水平的输出，而低质量的 Content 不仅无益，甚至可能引入错误的信息或约束，损害输出质量。

理解这一区分的关键在于：Description 的好坏决定了 Skill 能否\textit{被找到}，Content 的好坏决定了 Skill 被找到后能否\textit{发挥作用}。这两个维度是相互独立的。一个 Description 写得极为精准的 Skill，其 Content 可能是简单复制自其他 Skill 而毫无价值；反之，一个 Content 包含深度领域专业知识的 Skill，其 Description 可能过于笼统，在大量候选中难以被精确识别。现有平台的根本问题在于，它们以 Star 数或下载量作为统一的质量指标，而这两个指标对于 Description 质量和 Content 质量的区分度都极为有限——一个 Skill 的高 Star 数可能仅仅反映其 Description 写得吸引人，而非其 Content 确实有用。

以语义地形（Semantic Terrain）作为直观类比：若将每个 Skill 的 Description embedding 视为其在需求语义空间中的二维坐标（XY 位置），则功能相近的 Skill 自然聚集在地形图的同一区域，形成语义山峰。在同一个山峰内（即功能定位相近的 Skill 集合中），Content 的质量决定了每个 Skill 的"海拔高度"（Z 轴位置）。面向"Python 代码优化"的 Skill 与面向"浏览器自动化"的 Skill 分属两座不同的山峰，比较哪座山峰更高并无意义；有意义的比较是在同一座山峰内，找出海拔最高的那几个 Skill，即在功能定位相同的竞争者中找出 Content 质量最优的。

\begin{figure}[htbp]
\centering
\resizebox{0.9\textwidth}{!}{%
\begin{tikzpicture}[scale=0.85]
  % 绘制三维坐标轴
  \draw[->] (0,0,0) -- (6,0,0) node[right] {$X$ (Description 语义坐标 1)};
  \draw[->] (0,0,0) -- (0,5,0) node[above] {$Z$ (SkillRank 分值)};
  \draw[->] (0,0,0) -- (0,0,5) node[below left] {$Y$ (Description 语义坐标 2)};

  % 聚类 A：浏览器操作类 Skill（左侧山峰）
  \fill[blue!30,opacity=0.6] (1,0,1) -- (2.5,0,1) -- (2.5,0,2.5) -- (1,0,2.5) -- cycle;
  \fill[blue!50,opacity=0.5] (1.5,2.8,1.5) -- (2,2.8,1.5) -- (2,2.8,2) -- (1.5,2.8,2) -- cycle;
  \draw[blue!70] (1,0,1) -- (1.5,2.8,1.5);
  \draw[blue!70] (2.5,0,1) -- (2,2.8,1.5);
  \draw[blue!70] (2.5,0,2.5) -- (2,2.8,2);
  \draw[blue!70] (1,0,2.5) -- (1.5,2.8,2);
  % 高质量 Skill 标注
  \filldraw[red] (1.75,2.8,1.75) circle (3pt) node[above right, font=\small] {高质量 Skill};
  % 低质量 Skill
  \filldraw[gray] (1.2,1.0,1.2) circle (2pt);
  \filldraw[gray] (2.1,0.6,2.0) circle (2pt);
  \filldraw[gray] (1.8,1.5,2.2) circle (2pt);
  \node[blue!70, font=\small] at (1.75,-0.3,1.75) {浏览器操作聚类};

  % 聚类 B：代码生成类 Skill（右侧山峰）
  \fill[green!30,opacity=0.6] (3.5,0,1) -- (5.5,0,1) -- (5.5,0,3) -- (3.5,0,3) -- cycle;
  \fill[green!50,opacity=0.5] (4.0,3.5,1.5) -- (5.0,3.5,1.5) -- (5.0,3.5,2.5) -- (4.0,3.5,2.5) -- cycle;
  \draw[green!70] (3.5,0,1) -- (4.0,3.5,1.5);
  \draw[green!70] (5.5,0,1) -- (5.0,3.5,1.5);
  \draw[green!70] (5.5,0,3) -- (5.0,3.5,2.5);
  \draw[green!70] (3.5,0,3) -- (4.0,3.5,2.5);
  \filldraw[red] (4.5,3.5,2.0) circle (3pt) node[above right, font=\small] {高质量 Skill};
  \filldraw[gray] (4.0,1.2,1.3) circle (2pt);
  \filldraw[gray] (5.0,0.8,2.5) circle (2pt);
  \filldraw[gray] (4.3,2.0,1.8) circle (2pt);
  \node[green!70, font=\small] at (4.5,-0.3,2.0) {代码生成聚类};

  % 说明文字
  \node[font=\footnotesize, align=center] at (3,-0.8,3)
    {不同聚类间 Z 值不可比；\\同聚类内 Z 值排序即为 SkillRank};
\end{tikzpicture}
}
\caption{Skill 语义地形模型示意图。XY 平面坐标由 Description 向量降维得到，Z 轴高度为 SkillRank 分值。两个语义聚类代表不同功能领域，红点表示各聚类内的高质量 Skill。}
\label{fig:terrain}
\end{figure}

\section{Skill 价值评估的条件排序定义}
\label{sec:conditional_ranking}

基于上节的结构分析，本文给出 Skill 价值评估问题的精确定义。

\textbf{定义 2.1（语义邻域）} 设全量 Skill 集合为 $S = \{s_1, s_2, \ldots, s_n\}$，每个 Skill 的 Description 通过 Sentence-BERT 模型\cite{reimers2019sentencebert}编码为向量 $\mathbf{d}_i \in \mathbb{R}^{384}$。给定查询 context $q$（同样编码为向量 $\mathbf{e}_q \in \mathbb{R}^{384}$），查询 $q$ 对应的语义邻域 $\mathcal{N}(q)$ 定义为：

\begin{equation}
    \label{eq:neighborhood}
    \mathcal{N}(q) = \left\{ s_i \in S \;\middle|\; \cos(\mathbf{d}_i, \mathbf{e}_q) \geq \theta \right\}
\end{equation}

其中 $\theta$ 为相似度阈值，实践中通过检索 top-$M$ 最近邻（$M = \alpha K$，$\alpha \in [5, 10]$）来动态确定。

\textbf{定义 2.2（条件排序）} Skill 价值评估问题是一个条件排序问题：给定语义邻域 $\mathcal{N}(q)$ 和参数 $K$，在 $\mathcal{N}(q)$ 内部对所有 Skill 按 Content 质量排序，输出质量最高的前 $K$ 个 Skill 集合 $\mathcal{R}(q, K) \subseteq \mathcal{N}(q)$，$|\mathcal{R}(q, K)| = K$。

这一定义的"条件"性体现在排序仅在 $\mathcal{N}(q)$ 内部进行，不同语义邻域之间的 Skill 质量分不具备可比性。这与现有平台的全局排序方式存在根本区别：全局排序相当于对所有 $s_i \in S$ 而非仅对 $s_i \in \mathcal{N}(q)$ 进行排名，忽略了 Skill 价值的情境依赖特性。

此外，本文特别指出优化目标 $\mathcal{R}(q, K)$ 是一个\textbf{集合}而非排名第一的单个 Skill，这与 Agent 消费 Skill 的方式直接对应。Liu 等人的研究\cite{liu2024lostinthemiddle}表明，LLM 对上下文中不同位置的信息利用程度相对均匀，并不像人类那样对第一条结果有压倒性的注意力偏向。这意味着，在 $\mathcal{R}(q, K)$ 中第 $K$ 名的 Skill 对 Agent 输出质量的贡献，与第 1 名相比差距有限，因此集合层面的整体质量比个体的精确排序更为重要。

\section{SkillRank：基于 Content 亲缘关系的价值传播}
\label{sec:skillrank}

\subsection{从超链接图到 Content 亲缘图}

PageRank 算法\cite{brin1998anatomy}将网页的重要性定义为其在链接图上的平稳分布权重。其直觉理解是一个"随机冲浪者"模型：想象一个用户随机浏览网页，每到一个页面就随机点击其中某个链接跳转至下一页面，以概率 $1-d$ 随机跳至任意页面（对应阻尼因子 $d$ 的引入）。在这个模型中，被访问概率最高的页面，即为 PageRank 得分最高的页面。被越多高权重页面链接的页面，吸引随机冲浪者到达的概率越高，自身的 PageRank 也越高——这是一个自洽的递归价值定义。

将这一思想迁移至 Skill 价值评估，关键的障碍是：Skill 之间不存在显式的超链接关系。公开 GitHub Skill 仓库中的大量 Skill 往往是相互独立发布的，没有任何一个 Skill 主动"链接"至另一个 Skill。

本文的核心观察是：\textbf{Content 的语义相似性可以作为显式链接的结构等价物}。若 Skill $s_j$ 的 Content 与 Skill $s_i$ 的 Content 在语义空间中高度接近，可以认为 $s_j$ 的内容在某种程度上"参照"或"派生"自 $s_i$ 所代表的解法模式——无论这种参照是作者有意识地借鉴，还是在独立创作过程中不约而同地汇聚至同一高质量解法。当大量 Skill 的 Content 都与某个 Skill 的 Content 语义相似时，该 Skill 所代表的解法模式已经成为该语义领域内的"范式"，其质量得到了生态自发的广泛认可。这与 PageRank 的逻辑完全对称：被大量高质量网页链接的网页是重要的；被大量 Skill Content 高度相似引用的 Skill，是该语义邻域内质量最高的。

\subsection{SkillRank 的形式化定义}

设每个 Skill $s_i$ 的 Content 通过 Sentence-BERT 编码为向量 $\mathbf{v}_i \in \mathbb{R}^{384}$。利用 FAISS 近似最近邻搜索，为每个 Skill 检索 Content 语义空间中最近的 $K$ 个邻居（$K=15$），构建稀疏有向亲缘图 $G = (S, E)$，其中：

\begin{equation}
    \label{eq:affinity_graph}
    (s_i \to s_j) \in E \iff s_j \in \mathrm{topK\text{-}NN}_{\mathrm{Content}}(s_i),\quad i \neq j
\end{equation}

每条有向边 $(s_i \to s_j)$ 的权重为两端 Content 向量的余弦相似度：

\begin{equation}
    \label{eq:edge_weight}
    w_{ij} = \cos(\mathbf{v}_i, \mathbf{v}_j) = \frac{\mathbf{v}_i \cdot \mathbf{v}_j}{\|\mathbf{v}_i\| \cdot \|\mathbf{v}_j\|}
\end{equation}

在此稀疏加权图上，Skill $s_i$ 的 SkillRank 分值 $SR(s_i)$ 通过以下迭代方程求解：

\begin{equation}
    \label{eq:skillrank}
    SR(s_i) = (1 - d) + d \sum_{\substack{s_j \in S \\ (s_j \to s_i) \in E}} \frac{w_{ji}}{\sum_{s_k : (s_j \to s_k) \in E} w_{jk}} \cdot SR(s_j)
\end{equation}

其中 $d$ 为阻尼因子（经 AutoResearch 优化后取 $d = 0.82$）。方程的直觉是：Skill $s_i$ 的质量分由所有指向它的 Skill 贡献，贡献量与贡献者 $s_j$ 自身的质量分及其对 $s_i$ 的相似度权重成正比。通过幂迭代法（power iteration）求解，初始化 $SR_0(s_i) = 1/n$，迭代至收敛（$\|SR^{(t+1)} - SR^{(t)}\|_1 < 10^{-6}$）。

\begin{figure}[htbp]
\centering
\begin{tikzpicture}[
  node distance=2.2cm,
  skill/.style={circle, draw, fill=blue!15, minimum size=1cm, font=\small},
  highskill/.style={circle, draw=red, fill=red!20, minimum size=1.2cm, font=\small, line width=1.5pt},
  arrow/.style={->, thick, >=stealth}
]
  % 节点
  \node[highskill] (A) at (0, 0) {\shortstack{$s_1$\\\tiny SR=0.91}};
  \node[skill] (B) at (3, 1.5) {\shortstack{$s_2$\\\tiny SR=0.63}};
  \node[skill] (C) at (3, -1.5) {\shortstack{$s_3$\\\tiny SR=0.57}};
  \node[skill] (D) at (6, 2.5) {\shortstack{$s_4$\\\tiny SR=0.41}};
  \node[skill] (E) at (6, 0.5) {\shortstack{$s_5$\\\tiny SR=0.38}};
  \node[skill] (F) at (6, -1.5) {\shortstack{$s_6$\\\tiny SR=0.35}};
  \node[skill] (G) at (-3, 1) {\shortstack{$s_7$\\\tiny SR=0.44}};
  \node[skill] (H) at (-3, -1) {\shortstack{$s_8$\\\tiny SR=0.39}};

  % 边（指向 s1 表示相似）
  \draw[arrow, blue!60] (B) -- node[above, font=\tiny]{0.87} (A);
  \draw[arrow, blue!60] (C) -- node[below, font=\tiny]{0.84} (A);
  \draw[arrow, blue!60] (G) -- node[above, font=\tiny]{0.79} (A);
  \draw[arrow, blue!60] (H) -- node[below, font=\tiny]{0.76} (A);
  \draw[arrow, gray!60] (D) -- (B);
  \draw[arrow, gray!60] (E) -- (B);
  \draw[arrow, gray!60] (F) -- (C);
  \draw[arrow, gray!60] (D) -- (C);

  % 说明
  \node[font=\footnotesize, align=center, text width=5cm] at (1.5,-3.2)
    {$s_1$（红色节点）被4个邻居指向，\\Content被广泛相似引用，SkillRank最高};
\end{tikzpicture}
\caption{SkillRank 亲缘图示意（局部）。箭头表示 Content 相似关系，箭头上数字为余弦相似度权重。被大量节点以高权重指向的 $s_1$，获得最高的 SkillRank 分值（0.91）。}
\label{fig:skillrank_graph}
\end{figure}

\subsection{可扩展性：稠密图的稀疏化}

PageRank 得以在十亿级网页上运行，根本原因是网页超链接图天然稀疏——每个网页平均链接数约为数十，总边数与节点数呈线性关系。Skill 的 Content 相似度图则相反：任意两个 Skill 之间都存在非零的语义相似度，在 $n=10000$ 的规模下，全连接图的边数约为 $5 \times 10^7$，不仅存储代价不可接受，PageRank 的每轮迭代复杂度也将退化至 $O(n^2)$。

本文通过引入 FAISS IVF（倒排文件）索引进行 K-NN 截断，人为引入图稀疏性。对每个 Skill 仅保留 Content 语义空间中最近的 $K=15$ 个邻居，将每个节点的出度固定为 $K$，总边数降至 $O(nK)$。FAISS IVF 将向量空间划分为 $\lceil\sqrt{n}\rceil$ 个 Voronoi 单元，查询时仅探测候选单元（\texttt{nprobe} 参数控制），单次查询复杂度约为 $O(\log n)$，批量建图的整体复杂度为 $O(n\log n)$。在此稀疏图上，PageRank 每轮迭代复杂度为 $O(nK)$，$K$ 为常数，故迭代总复杂度为 $O(nK \cdot T)$（$T$ 为收敛轮数，实验中约 20 轮），整体线性于 $n$。

表 \ref{tab:complexity} 对比了 SkillRank 与朴素 $O(n^2)$ 全连接方案在理论复杂度上的差异。

\begin{table}[htbp]
\centering
\caption{SkillRank 与朴素全连接方案的复杂度对比}
\label{tab:complexity}
\small
\begin{tabular}{p{2.7cm}p{2.5cm}p{2.3cm}p{3.0cm}}
\toprule
\makecell[c]{方案} &
\makecell[c]{图构建\\复杂度} &
\makecell[c]{单轮迭代\\复杂度} &
\makecell[c]{增量更新\\复杂度} \\
\midrule
\makecell[l]{朴素全连接} &
\makecell[c]{$O(n^2 d)$} &
\makecell[c]{$O(n^2)$} &
\makecell[c]{$O(n)$\\（重算全图）} \\
\makecell[l]{SkillRank\\（FAISS K-NN）} &
\makecell[c]{$O(nd + n\log n)$} &
\makecell[c]{$O(nK)$} &
\makecell[c]{$O(K\log n)$\\（局部更新）} \\
\bottomrule
\end{tabular}
\end{table}

\section{本章小结}
\label{sec:chap2_summary}

本章建立了本文分析体系的理论基础。首先将 Skill 形式化为 $\langle Description, Content \rangle$ 二元结构，明确区分了语义定位维度（Description）与内容质量维度（Content）；其次提出了条件排序问题定义（定义 2.1 和 2.2），通过语义邻域的概念将排序问题限定在有意义的比较范围内；进而推导了 SkillRank 算法的数学形式（公式 \ref{eq:skillrank}），将 PageRank 的随机游走价值传播机制迁移至 Content 亲缘图；最后分析了通过 FAISS K-NN 截断实现线性可扩展性的工程方案（表 \ref{tab:complexity}）。这些分析为第三章的系统设计和第五章的实验验证提供了完整的方法论基础。
